% ============================================================
%  Metódy inžinierskej práce — Cvičenie 6, úloha 3 (prez-themes.tex)
%
%  Šablóna na experimentovanie s témou (\usetheme) a farebnosťou
%  (\usecolortheme) prezentácie. Samotný súbor sa neprekladá priamo:
%  riadok so značkou nižšie nahrádza skript build-themes.sh
%  konkrétnou témou a farebnosťou; ten potom vygeneruje a preloží
%  jednotlivé varianty do priečinka themes/.
%
%  Spustenie:  ./build-themes.sh
% ============================================================
\documentclass{beamer}

@@THMESETTINGS@@

\setbeamercovered{transparent}

\usepackage[T1]{fontenc}
\usepackage[utf8]{inputenc}
\usepackage{lmodern}
\usepackage[slovak,shorthands=off]{babel}

\title[Ukážka témy]{Ukážka témy a farebnosti beamer}
\subtitle{\vspace{2mm}Metódy inžinierskej práce 2026/27}
\author[Matvii Nosenko a kol.]{Matvii Nosenko, Volodymyr Nedoluha\\[1mm]Danylo Nikolenko, Rostyslav Mykytiuk}
\institute[FIIT STU]{Fakulta informatiky a informačných technológií\\Slovenská technická univerzita v Bratislave}
\date{\footnotesize \today}

\begin{document}

\begin{frame}[fragile=singleslide]
\titlepage
\end{frame}

\begin{frame}[fragile=singleslide]\frametitle{Prehľad}
\tableofcontents
\end{frame}

\section{Príznaky}
\begin{frame}[fragile=singleslide]\frametitle{Odrážky 1. a 2. úrovne}
\begin{itemize}
\item Prvá úroveň odrážok
	\begin{itemize}
	\item Druhá úroveň
	\item Ďalšia druhá úroveň
	\end{itemize}
\item Pokračovanie prvej úrovne
\end{itemize}
\end{frame}

\begin{frame}[fragile=singleslide]\frametitle{Ukážka bloku a textu}
Vplyv na čitateľnosť: nadpisy, päty a farebnosť blokov.
\begin{block}{Blok}
Text v \texttt{block} rámiku.
\end{block}
\begin{alertblock}{Upozornenie}
Zdôraznený text v \texttt{alertblock}.
\end{alertblock}
\end{frame}

\end{document}
